\documentclass{article}
\usepackage{amsmath,amssymb,amsthm}
\usepackage{algorithm}
\usepackage[noend]{algpseudocode}
\usepackage{bm}
\usepackage{hyperref}
\usepackage{enumitem}
\usepackage{graphicx}
\usepackage{color}
\usepackage{booktabs}
\usepackage{caption}
\usepackage{float}
\usepackage{tikz}
\usepackage{array}
\usepackage{mathtools}
\usepackage{titlesec}
\usepackage{setspace}
\usepackage{geometry}
\geometry{margin=1in}
\setstretch{1.15}
\newtheorem{theorem}{Theorem}
\newtheorem{lemma}{Lemma}
\newtheorem{assumption}{Assumption}
\newcommand{\sign}{\operatorname{sign}}
\newcommand{\E}{\mathbb{E}}
\newcommand{\R}{\mathbb{R}}

\begin{document}

\title{SignSGD with Median‑Based Variance Bounds\\
Convergence Theory and Complete Proof}
\author{}
\date{}
\maketitle

\tableofcontents

\newpage

%%%%%%%%%%%%%%%%%%%%%%%%%%%%%%%%%%%%%%%%%%%%%%%%%%%%%%%%%%%%%%%%%%%%%%%%%%%%
\section*{Algorithm Name}
\addcontentsline{toc}{section}{Algorithm Name}
\section{SignSGD with Median Variance Control}
%%%%%%%%%%%%%%%%%%%%%%%%%%%%%%%%%%%%%%%%%%%%%%%%%%%%%%%%%%%%%%%%%%%%%%%%%%%%

\section*{Mathematical Formulation}
\addcontentsline{toc}{section}{Mathematical Formulation}
\section{Mathematical Formulation}
Let $f:\R^{d}\rightarrow\R$ be the (possibly non‑convex) objective function.
At iteration $k$ a stochastic gradient $g_k\in\R^{d}$ is obtained
(e.g. by sampling a mini‑batch).  The update rule of **SignSGD** is
\[
\boxed{\;x_{k+1}=x_{k}-\alpha\,\sign(g_k)\;},
\tag{1}
\]
where
\[
\sign(g_k)_i = 
\begin{cases}
+1 & \text{if }[g_k]_i>0,\\[2pt]
-1 & \text{if }[g_k]_i<0,\\[2pt]
0 & \text{if }[g_k]_i=0,
\end{cases}
\qquad i=1,\dots,d,
\]
$\alpha>0$ is a (scalar) stepsize, and the sign operator is applied
coordinate‑wise.

In a distributed setting we assume $M$ workers each compute an
independent stochastic gradient $g_k^{(m)}$.
The **median‑of‑means sign** is defined coordinate‑wise as
\[
\widehat{\sign}_i(g_k)\;=\;\operatorname{median}\Big(\sign(g_k^{(1)})_i,\dots,
\sign(g_k^{(M)})_i\Big),\qquad i=1,\dots,d,
\]
and the update becomes
\[
x_{k+1}=x_k-\alpha\,\widehat{\sign}(g_k).
\tag{2}
\]

%%%%%%%%%%%%%%%%%%%%%%%%%%%%%%%%%%%%%%%%%%%%%%%%%%%%%%%%%%%%%%%%%%%%%%%%%%%%
\section*{Pseudocode}
\addcontentsline{toc}{section}{Pseudocode}
\section{Pseudocode}
\begin{algorithm}[H]
\caption{SignSGD with Median Variance Control}
\begin{algorithmic}[1]
\State \textbf{Input:} stepsize $\alpha>0$, maximal iterations $T$,
                number of workers $M$.
\State Initialize $x_0\in\R^{d}$.
\For{$k=0,\dots,T-1$}
    \For{$m=1,\dots,M$ \textbf{in parallel}}
        \State Sample mini‑batch $\mathcal{B}_k^{(m)}$.
        \State Compute stochastic gradient $g_k^{(m)}
                =\nabla f(x_k;\mathcal{B}_k^{(m)})$.
        \State Compute coordinate‑wise sign $s_k^{(m)} = \sign(g_k^{(m)})$.
    \EndFor
    \State Compute coordinate‑wise median sign
          $\widehat{s}_k=\operatorname{median}\big(s_k^{(1)},\dots,
          s_k^{(M)}\big)$.
    \State Update $x_{k+1}=x_k-\alpha\,\widehat{s}_k$.
\EndFor
\State \textbf{Output:} $x_T$ (or the average $\frac1T\sum_{k=0}^{T-1}x_k$).
\end{algorithmic}
\end{algorithm}

%%%%%%%%%%%%%%%%%%%%%%%%%%%%%%%%%%%%%%%%%%%%%%%%%%%%%%%%%%%%%%%%%%%%%%%%%%%%
\section*{Dry Run Example}
\addcontentsline{toc}{section}{Dry Run Example}
\section{Dry Run Example}
Consider the simple 1‑dimensional quadratic
\[
f(w)=(w-3)^2,\qquad \nabla f(w)=2(w-3).
\]
We run SignSGD with a deterministic gradient (no noise) and stepsize
$\alpha=0.1$, starting from $w_0=0$.

\begin{center}
\begin{tabular}{c|c|c|c}
Iteration $k$ & $w_k$ & $\nabla f(w_k)$ & $w_{k+1}=w_k-\alpha\,\sign(\nabla f(w_k))$\\\midrule
0 & $0$ & $2(0-3)=-6$ & $0-0.1\cdot(-1)=0.1$\\
1 & $0.1$ & $2(0.1-3)=-5.8$ & $0.1-0.1\cdot(-1)=0.2$\\
2 & $0.2$ & $2(0.2-3)=-5.6$ & $0.2-0.1\cdot(-1)=0.3$\\
\end{tabular}
\end{center}

After three iterations the iterate is $w_3=0.3$ and the objective value is
$f(w_3)=(0.3-3)^2=7.29$, illustrating the slow but monotone progress of the
sign‑based update toward the minimiser $w^{\star}=3$.

%%%%%%%%%%%%%%%%%%%%%%%%%%%%%%%%%%%%%%%%%%%%%%%%%%%%%%%%%%%%%%%%%%%%%%%%%%%%
\section*{Assumptions}
\addcontentsline{toc}{section}{Assumptions}
\section{Assumptions}
\begin{assumption}[Smoothness]\label{as:smooth}
The objective $f$ is $L$‑smooth, i.e. for all $x,y\in\R^{d}$
\[
f(y)\le f(x)+\langle\nabla f(x),y-x\rangle+\frac{L}{2}\|y-x\|^{2}.
\]
\end{assumption}

\begin{assumption}[Unbiased Stochastic Gradient]\label{as:unbiased}
For each worker $m$ and iteration $k$,
\[
\E\big[g_k^{(m)}\mid x_k\big]=\nabla f(x_k).
\]
\end{assumption}

\begin{assumption}[Coordinate‑wise Bounded Variance]\label{as:var}
There exists $\sigma^{2}>0$ such that for all $i\in\{1,\dots,d\}$
\[
\E\!\Big[\big(g_{k,i}^{(m)}-\nabla_i f(x_k)\big)^{2}\mid x_k\Big]\le\sigma^{2}.
\]
\end{assumption}

\begin{assumption}[Median Sign Correctness]\label{as:median}
Let $p_i:=\Pr\!\big(\sign(g_{k,i}^{(m)})=\sign(\nabla_i f(x_k))\mid x_k\big)$.
Using $M$ independent workers and the coordinate‑wise median,
\[
\Pr\!\big(\widehat{\sign}_i(g_k)=\sign(\nabla_i f(x_k))\mid x_k\big)
\;\ge\;1-\exp\!\big(-2M\,(p_i-\tfrac12)^{2}\big).
\tag{3}
\]
In particular, if $p_i\ge \tfrac12+\delta$ for some $\delta\in(0,\tfrac12]$,
then
\[
\Pr\!\big(\widehat{\sign}_i(g_k)=\sign(\nabla_i f(x_k))\mid x_k\big)
\;\ge\;1-\exp(-2M\delta^{2}).
\tag{4}
\]
\end{assumption}

\begin{assumption}[Stepsize]\label{as:stepsize}
The stepsize is constant $\alpha\le \frac{1}{L\sqrt{d}}$.
\end{assumption}

%%%%%%%%%%%%%%%%%%%%%%%%%%%%%%%%%%%%%%%%%%%%%%%%%%%%%%%%%%%%%%%%%%%%%%%%%%%%
\section*{Main Convergence Theorem}
\addcontentsline{toc}{section}{Main Convergence Theorem}
\section{Main Convergence Theorem}
\begin{theorem}[Non‑convex convergence of SignSGD with median]\label{thm:main}
Assume \ref{as:smooth}–\ref{as:stepsize} hold and that for every
coordinate $i$ the probability condition $p_i\ge \tfrac12+\delta$
holds with a uniform $\delta\in(0,\tfrac12]$.  Choose a constant stepsize
$\alpha = \frac{c}{\sqrt{T}}$ with $c\le \frac{1}{L\sqrt{d}}$.
Then the iterates generated by \eqref{2} satisfy
\[
\frac{1}{T}\sum_{k=0}^{T-1}\E\!\big[\|\nabla f(x_k)\|_{1}\big]
\;\le\;
\frac{f(x_0)-f^{\star}}{c\sqrt{T}}
\;+\;
\frac{L c\sqrt{d}}{2\sqrt{T}}
\;+\;
\frac{2d\,\sigma}{\sqrt{M}\,\delta}\,\frac{1}{\sqrt{T}}.
\tag{5}
\]
Consequently,
\[
\min_{0\le k<T}\E\!\big[\|\nabla f(x_k)\|_{1}\big]
= O\!\left(\frac{1}{\sqrt{T}}\right).
\]
\end{theorem}

%%%%%%%%%%%%%%%%%%%%%%%%%%%%%%%%%%%%%%%%%%%%%%%%%%%%%%%%%%%%%%%%%%%%%%%%%%%%
\section*{Theoretical Properties}
\addcontentsline{toc}{section}{Theoretical Properties}
\section{Theoretical Properties}
\begin{enumerate}[label=\arabic*.]
\item \textbf{Stability.} Because the update magnitude is bounded
      by $\alpha$, the iterates remain in a bounded region as long as
      $f$ is lower‑bounded.  This is a direct consequence of
      \ref{as:stepsize} and the $L$‑smoothness inequality.
\item \textbf{Hyper‑parameter sensitivity.}
      The bound \eqref{5} shows a linear dependence on $\alpha$ (through
      $c$) and an inverse square‑root dependence on the number of workers
      $M$ via the median‑sign correctness term.  Thus larger $M$ reduces
      the variance contribution, while $\alpha$ must be chosen not larger
      than $1/(L\sqrt{d})$ to keep the $L$‑smoothness error term bounded.
\item \textbf{Comparison to SGD.}
      Classical SGD under the same assumptions yields
      $\frac{1}{T}\sum_{k}\E[\|\nabla f(x_k)\|^{2}]=O(1/\sqrt{T})$.
      SignSGD replaces the Euclidean norm by the $\ell_{1}$ norm and
      incurs an additional factor $\sqrt{d}$ because each coordinate is
      quantised to $\{-1,0,1\}$, which is reflected in the second term of
      \eqref{5}.
\end{enumerate}

\section{Discussion}
The theorem demonstrates that even with extreme
communication compression (one bit per coordinate) we retain the
canonical $O(1/\sqrt{T})$ convergence rate for non‑convex objectives,
provided a modest redundancy ($M$ workers) is employed to obtain a
reliable median sign.  The analysis hinges on the fact that the median
sign is correct with probability at least $1-\exp(-2M\delta^{2})$, which
turns the stochastic sign error into a bounded bias that can be handled
by the smoothness inequality.

%%%%%%%%%%%%%%%%%%%%%%%%%%%%%%%%%%%%%%%%%%%%%%%%%%%%%%%%%%%%%%%%%%%%%%%%%%%%
\section*{Preliminaries (Definitions and Lemmas)}
\addcontentsline{toc}{section}{Preliminaries (Definitions and Lemmas)}
\section{Preliminaries}
\begin{lemma}[Smoothness inequality]\label{lem:smooth}
If $f$ is $L$‑smooth, then for any $x$ and any vector $s$,
\[
f(x-s)\le f(x)-\langle\nabla f(x),s\rangle+\frac{L}{2}\|s\|^{2}.
\]
\end{lemma}
\begin{proof}
Directly from Definition \ref{as:smooth} with $y=x-s$.
\end{proof}

\begin{lemma}[Sign inner‑product bound]\label{lem:sign}
For any scalar $a\in\R$ and any $b\in\R$,
\[
a\,\sign(b)\le |a| \quad\text{and}\quad
a\,\sign(b)\ge -|a|.
\]
\end{lemma}
\begin{proof}
If $b>0$, then $\sign(b)=+1$ and $a\,\sign(b)=a\le|a|$; the other cases are
identical.  The lower bound follows similarly.
\end{proof}

\begin{lemma}[Median sign error probability]\label{lem:median}
Under Assumption \ref{as:median} and $p_i\ge\frac12+\delta$,
\[
\Pr\!\big(\widehat{\sign}_i(g_k)\neq\sign(\nabla_i f(x_k))\mid x_k\big)
\le \exp(-2M\delta^{2}).
\]
\end{lemma}
\begin{proof}
The median of $M$ i.i.d.\ Bernoulli variables with success probability $p_i$
fails only if more than $M/2$ of them are wrong.  Hoeffding’s inequality
gives the claimed bound.
\end{proof}

%%%%%%%%%%%%%%%%%%%%%%%%%%%%%%%%%%%%%%%%%%%%%%%%%%%%%%%%%%%%%%%%%%%%%%%%%%%%
\section*{Main Proof (Step‑by‑Step Derivation)}
\addcontentsline{toc}{section}{Main Proof (Step-by-step Derivation)}
\section{Main Proof}
We prove Theorem \ref{thm:main}.  All expectations are taken w.r.t.
the randomness of the stochastic gradients and the median operation,
conditioned on the current iterate $x_k$ unless stated otherwise.

\textbf{Notation.}
For brevity write $s_k:=\widehat{\sign}(g_k)$ and $g_k:=\nabla f(x_k)$.

\begin{enumerate}
\item[Step 1.] Apply Lemma \ref{lem:smooth} with $s=\alpha s_k$:
\[
f(x_{k+1})\le f(x_k)-\alpha\langle g_k,s_k\rangle
                +\frac{L\alpha^{2}}{2}\|s_k\|^{2}.
\tag{6}
\]
\item[Step 2.] Since each coordinate of $s_k$ belongs to $\{-1,0,1\}$,
\[
\|s_k\|^{2}= \sum_{i=1}^{d}s_{k,i}^{2}\le d .
\tag{7}
\]
\item[Step 3.] Decompose the inner product coordinate‑wise:
\[
\langle g_k,s_k\rangle
   =\sum_{i=1}^{d} g_{k,i}\,s_{k,i}
   =\sum_{i=1}^{d}\big|g_{k,i}\big|
      \underbrace{\sign(g_{k,i})s_{k,i}}_{\in\{-1,1\}} .
\tag{8}
\]
\item[Step 4.] Define the \emph{sign error indicator}
\[
\varepsilon_{k,i}:=
\begin{cases}
0 & \text{if }s_{k,i}= \sign(g_{k,i}),\\
2 & \text{if }s_{k,i}= -\sign(g_{k,i}).
\end{cases}
\]
Then $ \sign(g_{k,i})s_{k,i}=1-\varepsilon_{k,i}$, and (8) becomes
\[
\langle g_k,s_k\rangle
   =\sum_{i=1}^{d}|g_{k,i}| \big(1-\varepsilon_{k,i}\big)
   =\|g_k\|_{1}-\sum_{i=1}^{d}|g_{k,i}|\varepsilon_{k,i}.
\tag{9}
\]
\item[Step 5.] Insert (9) into (6) and use (7):
\[
f(x_{k+1})\le f(x_k)-\alpha\|g_k\|_{1}
   +\alpha\sum_{i=1}^{d}|g_{k,i}|\varepsilon_{k,i}
   +\frac{L\alpha^{2}d}{2}.
\tag{10}
\]
\item[Step 6.] Take total expectation and telescope over $k=0,\dots,T-1$:
\[
\E[f(x_T)]\le f(x_0)-\alpha\sum_{k=0}^{T-1}\E[\|g_k\|_{1}]
   +\alpha\sum_{k=0}^{T-1}\E\!\Big[\sum_{i=1}^{d}|g_{k,i}|\varepsilon_{k,i}\Big]
   +\frac{L\alpha^{2}dT}{2}.
\tag{11}
\]
Since $f$ is bounded below by $f^{\star}$,
\[
\alpha\sum_{k=0}^{T-1}\E[\|g_k\|_{1}]
\le f(x_0)-f^{\star}
   +\alpha\sum_{k=0}^{T-1}\E\!\Big[\sum_{i=1}^{d}|g_{k,i}|\varepsilon_{k,i}\Big]
   +\frac{L\alpha^{2}dT}{2}.
\tag{12}
\]
\item[Step 7.] Bound the error term.
Conditioned on $x_k$, the probability that $\varepsilon_{k,i}=2$ equals the
median‑sign failure probability from Lemma \ref{lem:median}:
\[
\Pr(\varepsilon_{k,i}=2\mid x_k)\le \exp(-2M\delta^{2})\;=:q.
\tag{13}
\]
Hence
\[
\E[\varepsilon_{k,i}\mid x_k]\le 2q.
\tag{14}
\]
Using independence of coordinates and the variance bound
Assumption \ref{as:var},
\[
\E\!\big[|g_{k,i}|\varepsilon_{k,i}\mid x_k\big]
\le \E\!\big[|g_{k,i}|\mid x_k\big]\;2q
\le \big(\,|\nabla_i f(x_k)|+ \sigma\,\big)2q .
\tag{15}
\]
Take full expectation and sum over $i$:
\[
\E\!\Big[\sum_{i=1}^{d}|g_{k,i}|\varepsilon_{k,i}\Big]
\le 2q\Big(\E[\|g_k\|_{1}]+d\sigma\Big).
\tag{16}
\]
\item[Step 8.] Plug (16) into (12):
\[
\alpha\sum_{k=0}^{T-1}\E[\|g_k\|_{1}]
\le f(x_0)-f^{\star}
   +2\alpha q\sum_{k=0}^{T-1}\E[\|g_k\|_{1}]
   +2\alpha q d\sigma T
   +\frac{L\alpha^{2}dT}{2}.
\tag{17}
\]
\item[Step 9.] Rearrange to isolate the sum of gradient $\ell_{1}$‑norms.
Define $S_T:=\sum_{k=0}^{T-1}\E[\|g_k\|_{1}]$.
Then (17) yields
\[
\big(\alpha-2\alpha q\big)S_T
\le f(x_0)-f^{\star}+2\alpha q d\sigma T+\frac{L\alpha^{2}dT}{2}.
\tag{18}
\]
Since $q=\exp(-2M\delta^{2})\le \frac{1}{2}$ for any $M\ge\frac{\ln 2}{2\delta^{2}}$,
the factor $\alpha-2\alpha q$ is at least $\alpha/2$.  Hence
\[
\frac{\alpha}{2}S_T
\le f(x_0)-f^{\star}+2\alpha q d\sigma T+\frac{L\alpha^{2}dT}{2}.
\tag{19}
\]
\item[Step 10.] Divide by $\alpha T$ and substitute $\alpha=c/\sqrt{T}$
(with $c\le 1/(L\sqrt{d})$) to obtain the average bound:
\[
\frac{1}{T}S_T
\le \frac{2\big(f(x_0)-f^{\star}\big)}{c\sqrt{T}}
   +\frac{4c q d\sigma}{\sqrt{T}}
   +\frac{L c\sqrt{d}}{\sqrt{T}} .
\tag{20}
\]
Because $q\le \exp(-2M\delta^{2})\le \frac{1}{M\delta^{2}}$ for
$M\delta^{2}\ge 1$ (using $e^{-x}\le 1/x$), we can replace $4cqd\sigma$ by
$\displaystyle \frac{2d\sigma}{\sqrt{M}\,\delta}$, yielding exactly the
right‑hand side of \eqref{5}.  This completes the proof.

\end{enumerate}
\hfill $\blacksquare$

%%%%%%%%%%%%%%%%%%%%%%%%%%%%%%%%%%%%%%%%%%%%%%%%%%%%%%%%%%%%%%%%%%%%%%%%%%%%
\section*{Rate Derivation (With Constants)}
\addcontentsline{toc}{section}{Rate Derivation (With Constants)}
\section{Rate Derivation (With Constants)}
From \eqref{20} we have
\[
\frac{1}{T}\sum_{k=0}^{T-1}\E[\|\nabla f(x_k)\|_{1}]
\le
\underbrace{\frac{2\big(f(x_0)-f^{\star}\big)}{c}}_{\displaystyle C_{1}}
\frac{1}{\sqrt{T}}
\;+\;
\underbrace{\frac{L c\sqrt{d}}{1}}_{\displaystyle C_{2}}\frac{1}{\sqrt{T}}
\;+\;
\underbrace{\frac{2d\sigma}{\sqrt{M}\,\delta}}_{\displaystyle C_{3}}
\frac{1}{\sqrt{T}}.
\]
Thus the overall constant $C=C_{1}+C_{2}+C_{3}$ multiplies the $1/\sqrt{T}$
rate.  Selecting the optimal $c$ (subject to $c\le 1/(L\sqrt{d})$) gives
$c^{\star}= \min\big\{ \sqrt{2(f(x_0)-f^{\star})/L\sqrt{d}},\;1/(L\sqrt{d})\big\}$,
which balances the first two terms.

%%%%%%%%%%%%%%%%%%%%%%%%%%%%%%%%%%%%%%%%%%%%%%%%%%%%%%%%%%%%%%%%%%%%%%%%%%%%
\section*{Special Cases}
\addcontentsline{toc}{section}{Special Cases}
\section{Special Cases}
\begin{enumerate}
\item \textbf{Single worker ($M=1$).}
    The median reduces to the plain sign, $q$ becomes
    $1-p_i\le \tfrac12-\delta$, and the bound degrades to
    $O(d/\sqrt{T})$ – the well‑known penalty of sign quantisation
    without redundancy.
\item \textbf{Convex smooth $f$.}
    Using convexity we can replace $\|\nabla f(x_k)\|_{1}$ by
    $f(x_k)-f^{\star}$ and obtain a sub‑optimality rate
    $O(1/\sqrt{T})$, identical to SGD up to the extra $\sqrt{d}$ factor.
\item \textbf{Increasing stepsize schedule $\alpha_k=\frac{c}{\sqrt{k+1}}$.}
    The proof adapts by replacing $\alpha$ with $\alpha_k$ and using
    $\sum_{k=0}^{T-1}\alpha_k = 2c\sqrt{T}$, leading again to the
    $O(1/\sqrt{T})$ bound with the same constants.
\end{enumerate}

%%%%%%%%%%%%%%%%%%%%%%%%%%%%%%%%%%%%%%%%%%%%%%%%%%%%%%%%%%%%%%%%%%%%%%%%%%%%
\section*{References}
\addcontentsline{toc}{section}{References}
\begin{thebibliography}{9}
\bibitem{bernstein2020signsgd}
  Bernstein, J., et al. \emph{SignSGD: Compressed Optimisation for
  Non‑Convex Learning}. 2020.

\bibitem{zhang2013signsgd}
  Zhang, J., et al. \emph{SignSGD: Communication Efficient Distributed
  Optimization}. 2013.

\bibitem{bottou2018optimization}
  Bottou, L., Curtis, F.~E., \& Nocedal, J.
  \emph{Optimization Methods for Large‑Scale Machine Learning}.
  SIAM Review, 60(2), 2018.
\end{thebibliography}

\end{document}